\worksheettitle
  {M06\(\star\). Диагонали и области}
  {\modulemark{M06\(\star\)} Усложнение}

\studentnameblock

\begin{notebox}[Выводим правило]
  Из выбранной вершины нельзя провести диагональ к ней самой и к двум
  соседним вершинам. Проверьте это на рисунках, прежде чем использовать
  общее правило.
\end{notebox}

\begin{practicebox}[1. Найдите закономерность]
  Заполните таблицу для диагоналей из одной вершины:
  \[
    \begin{array}{c|c|c|c|c}
      \text{число вершин }n & 4 & 5 & 6 & 8\\ \hline
      \text{диагоналей} & \answerblank[10mm] & \answerblank[10mm] &
        \answerblank[10mm] & \answerblank[10mm]
    \end{array}
  \]
  Сформулируйте правило через $n$: \answerblank[45mm].
\end{practicebox}

\quadrilateralregionsfigure

\begin{practicebox}[2. Прочитайте сложную схему]
  В четырёхугольнике $ABCD$ проведены $AC$ и $BM$, их пересечение --- $P$.
  Назовите минимальные области, на которые разделилась фигура.
  \writinglines{3}
\end{practicebox}

\begin{practicebox}[3. Проверьте границы]
  Для каждой названной области перечислите вершины подряд по её границе.
  Объясните, почему точку пересечения $P$ нельзя пропускать.
  \writinglines{3}
\end{practicebox}

\begin{checkpointbox}[Исследование]
  В выпуклом семиугольнике выберите вершину и проведите все диагонали из неё.
  Сколько диагоналей и сколько треугольных частей получилось? Обоснуйте оба
  ответа. \writinglines{4}
\end{checkpointbox}
